\documentclass[11pt,letterpaper]{article}

\usepackage[utf8]{inputenc}
\usepackage[T1]{fontenc}
\usepackage[margin=1in]{geometry}
\usepackage{tikz}
\usepackage{xcolor}
\usepackage{amssymb}
\usepackage{tabularx}
\usepackage{booktabs}
\usepackage{longtable}
\usepackage{enumitem}
\usepackage{hyperref}
\usepackage{titlesec}
\usepackage{fancyhdr}
\usepackage{parskip}

\definecolor{lux-red}{HTML}{8B1A1A}
\definecolor{lux-grey}{HTML}{4A4A4A}

\hypersetup{
  colorlinks=true,
  linkcolor=lux-red,
  citecolor=lux-grey,
  urlcolor=lux-red,
  pdftitle={Lux Industries Inc.\ --- IP Enforcement Memorandum},
  pdfauthor={Woo Bin --- Chief Technology Officer, Lux Industries Inc.},
  pdfsubject={Confidential board memorandum},
}

\titleformat{\section}{\normalfont\Large\bfseries\color{lux-grey}}{\thesection}{0.6em}{}
\titleformat{\subsection}{\normalfont\large\bfseries\color{lux-grey}}{\thesubsection}{0.6em}{}

\pagestyle{fancy}
\fancyhf{}
\fancyhead[L]{\small\color{lux-grey}\textsc{Confidential --- Board Materials --- Attorney-Client Privileged}}
\fancyhead[R]{\small\color{lux-grey}Lux Industries Inc.}
\fancyfoot[C]{\small\color{lux-grey}\thepage}
\renewcommand{\headrulewidth}{0.2pt}

\begin{document}

% ─────────────────────────────────────────────────────────────────
% Cover page
% ─────────────────────────────────────────────────────────────────
\pagestyle{empty}
\begin{tikzpicture}[remember picture,overlay]
  \fill[black] (current page.south west) rectangle (current page.north east);

  \node[anchor=north west] at ([xshift=1.4cm,yshift=-1.4cm]current page.north west) {%
    \begin{tikzpicture}[scale=0.5]
      \fill[white] (-0.7,0.55) -- (0.7,0.55) -- (0,-0.55) -- cycle;
    \end{tikzpicture}%
  };

  \node[anchor=center,text=white,font=\fontsize{30}{38}\selectfont\bfseries,text width=14cm,align=center]
    at ([yshift=3.2cm]current page.center) {LUX INDUSTRIES};

  \node[anchor=center,text=white!70,font=\fontsize{14}{18}\selectfont,text width=14cm,align=center]
    at ([yshift=1.6cm]current page.center)
    {IP Enforcement Memorandum};

  \draw[white!30,line width=0.5pt]
    ([yshift=0.5cm,xshift=-5cm]current page.center) --
    ([yshift=0.5cm,xshift=5cm]current page.center);

  \node[anchor=center,text=white!40,font=\fontsize{9}{11}\selectfont]
    at ([yshift=-0.8cm]current page.center)
    {Prepared for: Cyrus Pahlavi, President \& CEO};

  \node[anchor=center,text=white!40,font=\fontsize{9}{11}\selectfont]
    at ([yshift=-1.4cm]current page.center)
    {Hunter Dupont, Managing Partner --- Strategic Transactions};

  \node[anchor=center,text=white!40,font=\fontsize{9}{11}\selectfont]
    at ([yshift=-2.0cm]current page.center)
    {Alan Donohue, General Counsel};

  \node[anchor=center,text=white!40,font=\fontsize{9}{11}\selectfont]
    at ([yshift=-2.8cm]current page.center)
    {Prepared by: Woo Bin --- CTO, Lux Industries Inc.};

  \node[anchor=center,text=white!30,font=\fontsize{9}{11}\selectfont]
    at ([yshift=-3.4cm]current page.center)
    {2026-05-25};

  \node[anchor=south west,text=white!20,font=\fontsize{7}{9}\selectfont]
    at ([xshift=1.4cm,yshift=0.8cm]current page.south west)
    {Lux Industries Inc.};

  \node[anchor=south east,text=white!20,font=\fontsize{7}{9}\selectfont]
    at ([xshift=-1.4cm,yshift=0.8cm]current page.south east)
    {CONFIDENTIAL --- ATTORNEY-CLIENT PRIVILEGED};
\end{tikzpicture}
\clearpage
\pagestyle{fancy}
\setcounter{page}{1}

% ─────────────────────────────────────────────────────────────────
\begin{flushleft}
\textbf{\textsc{Memorandum}}\\[0.5em]
\begin{tabularx}{\textwidth}{lX}
\textbf{TO:}        & Cyrus Pahlavi, President \& Chief Executive Officer, Lux Industries Inc. \\
                    & Hunter Dupont, Managing Partner --- Strategic Transactions (deal lead) \\
                    & Alan Donohue, General Counsel (legal lead) \\
\textbf{CC:}        & Board of Directors, Lux Industries Inc. \\
\textbf{FROM:}      & Woo Bin --- Chief Technology Officer, Lux Industries Inc.; concurrently Vice President of Engineering, the Counterparty \\
\textbf{DATE:}      & 25 May 2026 \\
\textbf{RE:}        & Misappropriation of Lux foundational technology by the Counterparty (Satschel, Inc.\ d/b/a Liquidity.io, and operating affiliate Liquidity.io LLC); legal options, damages analysis, and proposed cross-license and settlement framework \\
\textbf{PRIVILEGE:} & Confidential. Attorney-client privileged. Attorney work product. Prepared in anticipation of litigation and as a board strategy document. \\
\end{tabularx}
\end{flushleft}

\vspace{0.3em}
\begin{small}
\textit{Author's disclosure of interest:} I serve concurrently as
Chief Technology Officer of Lux Industries Inc.\ and Vice President
of Engineering at the Counterparty entity referenced herein. This
memorandum is written from the Lux chair, in service of the Lux
board's decision-making. My Counterparty role gives direct
operational visibility into the matters described; I have set them
out as observed, with the understanding that any Counterparty-side
response will be developed by its own counsel and that Lux's outside
litigation counsel will need to develop independent record support
for any allegations herein before filing.
\end{small}

\vspace{0.5em}\hrule height 0.6pt\vspace{1em}

% ─────────────────────────────────────────────────────────────────
\section*{I.\ Executive Summary}

\textbf{The Counterparty has commercially deployed material portions
of Lux Industries Inc.'s patent-protected and source-available
proprietary technology since at least 2014, in production, on three
live blockchain networks, and in support of a registered ATS/BD/TA
business that is currently raising approximately \$75 million at a
\$1.1 billion valuation, with \textbf{Forum Markets} (a publicly-traded
company, formerly Ethzilla; ticker \textbf{[TICKER --- to be confirmed
against SEC EDGAR]}) continuing as lead investor from the prior round
and verbal commitments for at least approximately half the round
secured to date --- in each case without executing the commercial
license that the governing Lux Ecosystem License v1.2
(``LEL'') and the Lux Research License with Patent Reservation v1.0
(``LRL--PR'') require.}

\textbf{Lux's direct investment in the misappropriated technology
base exceeds \$20 million in development cost recorded since 2014,
plus a multi-year contribution from senior engineering personnel
including the Lux Chief Cryptographer and the named author of this
memorandum.}

\textbf{The most striking evidentiary fact in this memorandum is the
90-day commit-ratio analysis at \S VI:} the Lux Founder and Chief
Cryptographer contributed \textbf{5,991 commits} and approximately
\textbf{93 million lines of code} (additions + deletions) across six
venture trees --- the Counterparty's primary codebase plus five
related ventures (VCC, MLC, EquityTable, OnyxPlus, and the
legacy/archive codebase including the simplici-* repositories) ---
during the most recent 90-day window, and \textbf{zero commits to
Lux's upstream repositories} in the same window. The Counterparty
reached its \$1.1B current-round valuation during that same 90 days.
The implied value-creation rate of the Lux Founder's diverted
bandwidth is approximately \textbf{\$12.2M per day}.

\textbf{Discovery posture.} This matter was identified and elevated to
the Lux Industries Inc.\ board by Mr.\ Woo Bin (Lux CTO; concurrent
VPE at the Counterparty), whose dual-role visibility --- \texttt{luxfi/*}
GitHub-org admin on the Lux side combined with operational seat on the
Counterparty's engineering organization --- positioned him to observe
the unauthorized integration of Lux modules into the Counterparty's
commercial stack and the absence of a corresponding executed license.
Mr.\ Bin remains in his VPE role; his team's continued involvement is
required for the Counterparty's long-term operational success, and the
desired outcome of this matter is an \textbf{equitable settlement
preserving Counterparty operational continuity while recognising Lux's
contribution}, rather than scorched-earth litigation.

\textbf{This memorandum sets out the legal record, identifies twelve
causes of action available to Lux, quantifies damages at small /
medium / big scenarios (\$80M to \$590M for IP claims; \$180M to
\$890M with opportunity-cost component added), projects forward
enterprise value at user-scale milestones up to 100M users (Lux's
10\% equity stake value ranges \$5M -- \$40B), analyzes the
``wait vs.\ settle'' time-value question, presents a settlement
triangulation matrix the deal lead (Mr.\ Dupont) can use to
negotiate, and locks the mandatory terms the President \& CEO (Mr.\
Pahlavi) requires in any executed agreement.}

\textbf{This is not a single-proposal document.} It presents data,
damage ranges, partnership-upside ranges, and negotiating floors.
The deal team triangulates to a specific outcome. Mr.\ Pahlavi
retains final approval authority on any executed agreement.

% ─────────────────────────────────────────────────────────────────
\section*{II.\ Parties and Dramatis Personae}

\subsection*{Lux Industries Inc.\ --- Decision Chain}

\begin{itemize}[leftmargin=1.5em,itemsep=0.2em]
\item \textbf{Cyrus Pahlavi} --- President \& Chief Executive Officer.
Final approval authority on any settlement, license, or litigation
decision.
\item \textbf{Hunter Dupont} --- Largest investor in Lux Industries
Inc.; Managing Partner, Strategic Transactions. Leads deal discussions
and commercial negotiations with the Counterparty.
\item \textbf{Alan Donohue} --- General Counsel. Leads legal strategy
and definitive-document drafting; provides legal support to Mr.\
Dupont throughout the negotiation.
\item \textbf{Woo Bin} --- Chief Technology Officer of Lux Industries
Inc.; \texttt{luxfi/*} GitHub-org admin. Author of this memorandum.
Holds concurrent role as Vice President of Engineering at the
Counterparty (\textbf{role currently active}; his team is required
for the Counterparty's long-term operational success). It was through
this concurrent visibility that Mr.\ Bin identified the unauthorized
integration of Lux modules into the Counterparty's commercial stack
and elevated the matter to the Lux Industries Inc.\ board, occasioning
this memorandum.
\item \textbf{Zach Kelling} --- Founder, Architect, and Chief
Cryptographer of Lux Industries Inc.; \texttt{luxfi/*} GitHub-org
admin; minority shareholder. \textbf{Currently serves as Chief
Technology Officer of the Counterparty}, with operational authority
over the platform rebuild and ongoing technical posture; the
Counterparty's commercial decisions and investor-facing
representations are owned by Mr.\ Choi (CEO) and Mr.\ Tromblye (CSO).
Mr.\ Kelling's stated objective in supporting this matter is an
equitable settlement preserving Counterparty operational continuity.
Not currently a director on either side. \textbf{See \S IV.E for
particularized treatment and \S XII for protective treatment.}
\end{itemize}

\subsection*{Counterparty --- Officer Slate}

\begin{itemize}[leftmargin=1.5em,itemsep=0.2em]
\item \textbf{Eric Choi} --- Chief Executive Officer of the
Counterparty. Top operational authority; ultimate responsibility for
the Counterparty's representations to investors and regulators
regarding IP ownership and licensing.
\item \textbf{Austin Tromblye} --- Chief Strategy Officer; previously
Acting Chief Executive Officer. Blockchain / AI engineer. Built a
\textbf{prior} (not current) ATS infrastructure at a prior employer
(Franklin Templeton) and \textbf{served as an eyewitness in the
Franklin Templeton litigation} arising from that engagement.
\textbf{Counsel should review the Franklin Templeton litigation
record} for (i) pattern-of-conduct evidence relevant to the present
matter, (ii) any limits on or adverse characterizations of his prior
testimony that could be impeachable, and (iii) the technical scope and
quality of the ATS he built there.

Material context: prior to Mr.\ Kelling's arrival as active Chief
Technology Officer at the Counterparty, the Counterparty's ATS as
built under Mr.\ Tromblye's technical direction \textbf{lacked even
basic time-price priority order matching}, the foundational principle
of any functioning alternative trading system. The pre-Kelling
technical state is addressed in \S IV(F).
\item \textbf{Coleman Church} --- Chief Executive Officer of the
Counterparty's broker-dealer affiliate. FINRA-regulated responsibility
for the BD's technology and operational integrity.
\item \textbf{Jayant Sonsurkar} --- Chief Security Officer of the
Counterparty. Long-standing owner / admin on the Counterparty's GitHub
organizations: historically \texttt{satschel/*}, and currently
\texttt{liquidityio/*} --- the latter being the org under which the
rebuild integrating unlicensed Lux modules is being performed. The
CSO role at any regulated digital-asset venue carries the duty to
oversee software-supply-chain integrity, including open-source-license
review and third-party-IP compliance. As admin on the Counterparty's
own repositories, Mr.\ Sonsurkar had on-platform visibility into the
imported Lux LICENSE files --- LEL, LRL--PR, BSD-3, and others ---
as they entered the Counterparty's build tree. (For clarity:
\texttt{luxfi/*} GitHub-org admin was held by Mr.\ Bin and Mr.\
Kelling on the Lux side; Mr.\ Sonsurkar did not have, and was not
required to have, admin access on Lux's own repositories. His
duty of OSS-license review attaches to imports into the
Counterparty's tree, where his admin access was present.)
\end{itemize}

\subsection*{Counterparty Entity Structure}

For purposes of this memorandum, ``the Counterparty'' refers
collectively to the operating entities under the current capital
raise (a Delaware corporation parent and its US LLC operating
affiliate, transacting as Liquidity.io) unless context requires
otherwise. The corporate structure should be confirmed during
outside-counsel diligence.

\textbf{Lead investor (current round).} Forum Markets, a publicly-
traded company (formerly Ethzilla; ticker \textbf{[TICKER --- to be
confirmed against SEC EDGAR]}), continues in the lead-investor role
from the prior round and has signalled continued leadership of the
current \$75M raise at the \$1.1B valuation. Forum Markets' public-
issuer status means any decision to anchor a material investment in
the Counterparty is potentially material to Forum Markets' own
shareholders, with attendant Form 8-K considerations.

% ─────────────────────────────────────────────────────────────────
\section*{III.\ Jurisdiction, Venue, and Applicable Law}

\textbf{Federal causes of action} sound in: Defend Trade Secrets Act
(DTSA), 18 U.S.C.\ \S\S 1836--1839; Copyright Act, 17 U.S.C.\ \S\S
101 \emph{et seq.}; Patent Act, 35 U.S.C.\ \S\S 271, 281, 284, 285;
Lanham Act, 15 U.S.C.\ \S 1125(a); Securities Exchange Act of 1934,
15 U.S.C.\ \S 78j(b) (Rule 10b-5, 17 C.F.R.\ \S 240.10b-5).
\textbf{State and common-law claims} sound in Delaware (breach of
fiduciary duty, \emph{Aronson v.\ Lewis}, 473 A.2d 805 (Del.\ 1984);
aiding and abetting, \emph{Malpiede v.\ Townson}, 780 A.2d 1075
(Del.\ 2001)), and California (the likely Counterparty principal
place of business): common-law fraud, unjust enrichment, conversion,
tortious interference with prospective economic advantage, and civil
conspiracy.

\textbf{Forum recommendation:} Federal District Court for the
Northern District of California for the principal action, with
related Delaware Chancery proceeding for the corporate-governance
counts. Personal jurisdiction over the individual Counterparty
officers is established by their commercial activities directed at
Lux (a Delaware corporation) and their California residence.

\textbf{Litigation hold} --- Lux should issue a formal preservation
notice to the Counterparty and to each named individual within 30
days of any term-sheet delivery. Failure of preservation supports an
adverse-inference instruction at trial.

% ─────────────────────────────────────────────────────────────────
\section*{IV.\ Statement of Facts}

\textit{The following allegations are stated on information and
belief, based on the personal knowledge of the undersigned in his
dual operational role, supplemented by Lux's internal records.}

\subsection*{A. Background of Lux Industries Inc.\ and Its IP (2014--2026)}

\textbf{1.} Lux Industries Inc.\ (with continuous corporate identity
through Lux Partners Limited predecessor since 2014) has developed,
since 2014, a blockchain and cryptographic technology stack including
consensus algorithms, post-quantum cryptographic primitives, threshold
and MPC protocols, fully-homomorphic-encryption (FHE) precompile
architecture, virtual-machine substrates, EVM derivatives, decentralized
exchange matching engines, and cross-chain messaging protocols.

\textbf{2.} Lux's direct development investment in this stack exceeds
\$20 million in recorded cost from 2014 to date. Fully-loaded
personnel cost over the same period is a multiple of that figure.

\textbf{3.} Chronological development milestones:

\begin{itemize}[leftmargin=2em,itemsep=0.1em]
\item \textbf{2014 -- 2017}: Foundational consensus research; early
Quasar prototypes; cryptographic primitives library;
\texttt{secp256k1}, \texttt{Ed25519}, \texttt{BLS12-381} reference
implementations.
\item \textbf{2018 -- 2020}: Active development of the planet-scale
DEX matching engine with 1ms finality. Initial EVM fork (LGPL).
Threshold-MPC research.
\item \textbf{2020 -- 2023}: Quasar consensus stabilization; Pulsar
and Magnetar Apache-2.0 threshold-PQ primitives; FHE prototypes.
\item \textbf{2024}: Ringtail lattice-signature scheme; FIPS 203 /
204 / 205 verifier precompiles (ML-KEM, ML-DSA, SLH-DSA); P3Q
protocol; LEL v1.2 and LRL--PR v1.0 formalized; Warp messaging
protocol stabilization; FHE-based dark-pool matching.
\item \textbf{2025}: LX Suite DEX precompile family (Pool / Oracle /
Router / Hooks at LP-9010 onward) released following Uniswap v4's
hooks model; Corona Ring-LWE threshold scheme (successor to Ringtail);
LP-4200 precompile address standardization; multi-arch GPU acceleration
under the closed \texttt{\textasciitilde/work/luxcpp/*} tree.
\item \textbf{2026}: Brand-neutral cleanup; relicensing of Pulsar /
Magnetar under Apache-2.0; ongoing institutional partnership
development.
\end{itemize}

\textbf{4.} Lux holds or has filed for patents on, without limitation:
Quasar BFT consensus; sampler-bias correction; Ringtail / Corona
lattice threshold schemes; Unified VM Execution Interface with
Pluggable Backends; Post-Quantum State Verification with Verkle
Proofs; FHE precompile architecture; GPU-accelerated DEX matching
engine; AI-optimized bytecode execution.

\subsection*{B. The Lux Licensing Tier Structure}

\textbf{5.} Lux publishes its software under a three-tier strategy:

\textbf{6.} \textbf{Tier 1 (BSD-3-Clause)}: commodity infrastructure,
free commercial use. \texttt{luxfi/node}, \texttt{luxfi/sdk},
\texttt{luxfi/cli}, \texttt{luxfi/api}, \texttt{luxfi/codec},
\texttt{luxfi/database}, \texttt{luxfi/p2p}, and approximately twenty-
five (25) other modules.

\textbf{7.} \textbf{Tier 2 (LEL v1.2)}: source-available, patent-
protected. Per LEL \S\S 1 and 3, commercial rights granted only to
``Authorized Networks'' --- the Lux Primary Network (NetworkID 1,
ChainID 96369), Lux testnets and devnets, and ``L1/L2/L3 chains
descending from the Lux Primary Network.''

\textbf{8.} Per LEL \S 4, forks and competing networks prohibited.
Per LEL \S 5, all patent rights reserved. Per LEL \S 8, license
terminates automatically on breach. Modules: \texttt{luxfi/consensus},
\texttt{luxfi/crypto}, \texttt{luxfi/mpc}, \texttt{luxfi/threshold},
\texttt{luxfi/ringtail}, \texttt{luxfi/corona}, \texttt{luxfi/fhe},
\texttt{luxfi/precompile}, \texttt{luxfi/captable},
\texttt{luxfi/broker}, \texttt{luxfi/treasury},
\texttt{luxfi/transfer}, \texttt{luxfi/compliance},
\texttt{luxfi/hsm}, \texttt{luxfi/aml}, \texttt{luxfi/maker},
\texttt{luxfi/warp}, among others.

\textbf{9.} \textbf{Tier 3 (LRL--PR v1.0)}: research-only license.
Section 3 expressly prohibits commercial use; Section 4 reserves all
patent claims. Modules: \texttt{luxfi/vm}, \texttt{luxfi/staking},
\texttt{luxfi/sampler}, \texttt{luxfi/dex} (Go bindings prior to its
2026 move to fully proprietary).

\textbf{10.} \textbf{Tier 4 (Closed Proprietary)}: the
\texttt{\textasciitilde/work/luxcpp/*} tree (CUDA / Metal / WebGPU /
FPGA / lattice GPU accelerators, the C++ DEX matching engine). Not
currently in the Counterparty's production graph; access requires a
separately negotiated Private-Moat Addendum.

\subsection*{C. The Counterparty's Commercial Use of Lux IP}

\textbf{11.} The Counterparty operates a registered Alternative
Trading System, broker-dealer, and transfer-agent business in the
United States, commercially deploying the validator binary
\texttt{lqd} (built from \texttt{liquidityio/node}) on three live
blockchain networks (Devnet ChainID 8675312, Testnet ChainID 8675310,
Mainnet ChainID 8675309), each of which incorporates dozens of Lux
modules.

\textbf{12.} The Counterparty further deploys the Liquidity EVM,
ATS, BD, and TA backend services, white-label tenants (MLC, VCC,
TZERO, ReTokens, and others), and approximately 12,796 SecurityTokens
on Devnet alone running on top of Lux precompiles.

\textbf{13.} The companion audit (referenced internally as the
license inventory) catalogues every Lux module directly imported.
Approximately fifteen (15) directly-imported modules carry LEL or
LRL--PR tier licenses requiring an executed commercial license.

\textbf{14.} \textbf{No commercial license has ever been executed
between Lux Industries Inc.\ and the Counterparty.} Outside counsel
should verify the non-existence of any such instrument through
document discovery.

\textbf{15.} The Liquid chains do not currently activate the Lux
Warp messaging precompile (canonical address
\texttt{0x0200000000000000000000000000000000000005}). \textbf{Under
the LEL Authorized-Network definition, a chain that does not maintain
Lux consensus and Warp interop does not qualify for the LEL ecosystem-
permissive grant.} The Liquid chains, as currently configured, do
not qualify in code.

\subsection*{D. The Counterparty Officer Team and Knowledge}

\textbf{16.} Mr.\ Eric Choi (CEO) holds ultimate operational
responsibility for technology licensing and IP compliance. His
signature appears (or will appear) on the subscription documents for
the current \$75M / \$1.1B-valuation raise; representations therein
sound in his name. As of the date of this memorandum the round has
not closed; the duty of disclosure is therefore current and ongoing.

\textbf{17.} Mr.\ Austin Tromblye (CSO, prev.\ Acting CEO) ---
discussed in \S II above, including the Franklin Templeton litigation
history and the pre-Kelling ATS technical state.

\textbf{18.} Mr.\ Coleman Church (BD CEO) --- FINRA-regulated
responsibility for the BD's technology integrity; representations on
FINRA filings sound in his name.

\textbf{19.} Mr.\ Jayant Sonsurkar (CSO) --- Long-standing owner /
admin on the Counterparty's GitHub organizations \texttt{satschel/*}
(historic) and \texttt{liquidityio/*} (current rebuild org), giving
him direct, contemporaneous, ongoing visibility into every Lux
\texttt{LICENSE} file \textbf{as those files entered the
Counterparty's build tree} on import --- the LEL, LRL-PR, BSD-3,
GPL, LGPL, BUSL, MIT, and Apache markers travel with their source
files into any tree that consumes them. \texttt{luxfi/*} GitHub-org
admin was held by Lux's own engineering team (Mr.\ Bin, Mr.\ Kelling)
on the upstream side; Mr.\ Sonsurkar did not have admin on
\texttt{luxfi/*} itself, nor was he required to --- his
software-supply-chain integrity duty attaches to imports into the
Counterparty's own tree, where his admin access was present.

\textbf{20.} The CSO role at a regulated digital-asset venue
carrying ATS / BD / TA registrations includes (without limitation)
software-supply-chain integrity oversight, identification of LEL /
LRL--PR / GPL / LGPL restricted licenses in deployed dependencies,
escalation of licensing mismatches to the CEO / BD CEO / GC, and
implementation of remediation.

\textbf{21.} On information and belief, Mr.\ Sonsurkar did not, at
any time prior to the elevation of this matter by Mr.\ Bin, raise the
LEL / LRL--PR mismatch internally, request a commercial license,
include the unlicensed-IP status in regulatory compliance reviews, or
flag the issue during current-raise diligence. By contrast, Mr.\ Bin
(VPE at the Counterparty, concurrently CTO at Lux) did identify the
mismatch through his own dual-role visibility and elevated it
upstream. The omission by Mr.\ Sonsurkar --- the officer to whom that
identification properly belonged --- is the central evidentiary fact
supporting Count IV (breach of fiduciary duty) and Count V (aiding
and abetting).

\textbf{22.} Messrs.\ Choi, Tromblye, and Church each had operational
or strategic visibility into the Counterparty's technology stack
during the relevant period; the precise scope of each officer's
knowledge will require diligence development.

\subsection*{E. The Position of Mr.\ Zach Kelling}

\textbf{23.} Mr.\ Kelling is the Founder of both Lux Industries
Inc.\ and the Counterparty. He \textbf{currently serves as Chief
Technology Officer of the Counterparty}. His operational authority
covers the platform rebuild and the ongoing technical posture; the
Counterparty's \textbf{commercial decisions and investor-facing
representations are owned by Mr.\ Choi (CEO) and Mr.\ Tromblye (CSO)},
not by Mr.\ Kelling. The dual role with Lux Industries Inc.\ (Founder,
Architect, Chief Cryptographer) is contemporaneous and material to the
discovery and disposition of this matter.

\textbf{24.} Mr.\ Kelling is a \textbf{minority shareholder of Lux
Industries Inc.}\ Counterparty equity position to be confirmed in
diligence; the Counterparty's officer slate as represented to
investors should be reconciled with any equity-grant records as part
of the fact-development for any negotiation or filing.

\textbf{25.} The factual record supports a finding that Mr.\ Kelling
is \textbf{operationally engaged on technical rebuild and not engaged
in the Counterparty's commercial direction or investor-facing
representations}. His stated objective in supporting the elevation of
this matter is an \textbf{equitable resolution that preserves the
Counterparty's operational continuity and Mr.\ Bin's team's role in
that continuity}, while recognising Lux Industries Inc.'s underlying
contribution. He is not a principal actor in the commercial-
representation conduct that grounds Counts IX (Securities Fraud), XI
(Common-Law Fraud), or X (Lanham Act). \textbf{See \S XII for
protective treatment, narrowed to his stated operational scope.}

\subsection*{F. Pre-Kelling Technical State and the ACH Exploit}

\textbf{26.} In the period immediately preceding Mr.\ Kelling's
arrival as active Chief Technology Officer, the Counterparty's ATS
platform was materially non-functional as a regulated alternative
trading system. Specifically, on information and belief:

\begin{itemize}[leftmargin=2em,itemsep=0.1em]
\item The ATS \textbf{did not implement working time-price priority
order matching} --- the foundational ordering principle of every
regulated US alternative trading system, required by SEC Reg ATS
for any venue purporting to operate as a securities trading facility.
\item No functioning ATS pipeline existed end-to-end.
\item \textbf{One day prior to Mr.\ Kelling's arrival as active CTO,
the platform suffered an ACH-funding exploit that resulted in the
theft of approximately \$40,000} via a trivial attack vector that
should have been remediated by even baseline security review. The
exploit evidences both the security posture of the platform under
the pre-Kelling officer team and the absence of competent CSO
oversight.
\end{itemize}

\textbf{27.} Upon his arrival, Mr.\ Kelling and his team (substantially
overlapping with the Lux engineering team) rebuilt the ATS / BD / TA
stack from the ground up, importing the Lux modules described in
\S IV(C) as the foundation. \textbf{In substantial part, the
technology that supports the current \$75M / \$1.1B-valuation raise was the Lux
IP that Mr.\ Kelling brought across, integrated with the regulatory
wrapper.}

\textbf{28.} The factual proposition that the Counterparty's
enterprise value as represented to investors is materially attributable
to Lux-origin technology contributed by Mr.\ Kelling and the broader
Lux team is supported by the contemporaneous technical record. This
context bears on damages attribution (\S VI), Mr.\ Tromblye's
credibility as a technology witness, the breach-of-fiduciary-duty
count, and the protective treatment of Mr.\ Kelling (\S XII).

\subsection*{G. The Current \$75M / \$1.1B-Valuation Capital Raise (Forum Markets, Lead)}

\textbf{29.} The Counterparty is \textbf{currently undertaking a
capital raise of approximately \$75 million at a \$1.1 billion
valuation}. The round is led by \textbf{Forum Markets} (a
publicly-traded company, formerly Ethzilla; ticker \textbf{[TICKER ---
to be confirmed against SEC EDGAR]}), continuing in the lead-investor
role from the prior round. As of the date of this memorandum, the
Counterparty has \textbf{verbal commitments covering at least
approximately half of the \$75 million round} (i.e.\ approximately
\$37.5M soft-circled). \textbf{The round has not closed.} Funds have
not been wired against current-round subscriptions; subscribers have
not been issued securities in respect of the current round; the
Counterparty's current and ongoing duty of disclosure to those
subscribers remains live until closing.
The subscription agreements, PPM, and side letters require outside-
counsel review for representations regarding (a) Counterparty
ownership of the underlying technology stack, (b) licensing status of
all third-party IP, (c) related-party transactions, (d) the CSO's
and BD CEO's compliance attestations, and (e) technology disclosures
in the risk factors section.

\textbf{30.} To the extent the current-raise subscription documents represent or omit
material facts regarding the Counterparty's ownership or licensed
access to Lux IP, those representations are independently actionable
under Section 10(b) of the Securities Exchange Act and Rule 10b-5.
See Count IX, infra.

% ─────────────────────────────────────────────────────────────────
\section*{V.\ Causes of Action Available to Lux Industries Inc.}

\subsection*{Count I --- Misappropriation of Trade Secrets (DTSA, 18 U.S.C.\ \S 1836)}

Lux's unpublished optimizations, integration know-how, architectural
documents, deployment recipes, MPC operational procedures, sampler-
bias remediation, and pre-publication private-repository code qualify
as trade secrets. Mr.\ Sonsurkar's continuing repository access
without an executed license, coupled with the licensing-mismatch
silence, supports misappropriation by improper means.

\textbf{Remedies}: Injunction (18 U.S.C.\ \S 1836(b)(3)(A)); actual
damages and unjust enrichment (\S (B)); reasonable royalty in lieu
(id.); exemplary damages up to 2$\times$ for willful and malicious
misappropriation (\S (C)); attorneys' fees (\S (D)); seizure if
record warrants (\S 1836(b)(2)).

\subsection*{Count II --- Copyright Infringement (17 U.S.C.\ \S 501)}

Lux owns copyright in every source file under \texttt{luxfi/*}.
Copyright notices are present in source headers. The Counterparty
has copied by direct import in \texttt{go.mod}, container-image
embedding, binary distribution to validator nodes, and on-chain
deployment of derivative SecurityTokens.

\textbf{Remedies}: Actual damages and infringer's profits, \S 504(b);
or statutory damages up to \$150,000 per work willfully infringed,
\S 504(c)(2); injunction (\S 502); impoundment (\S 503); attorneys'
fees and costs (\S 505).

\subsection*{Count III --- Patent Infringement (35 U.S.C.\ \S 271)}

Lux holds or has filed for patents covering the inventions
enumerated in \S IV(A)(4). The Counterparty's deployment of Quasar
consensus, Ringtail / Corona thresholds, the sampler-bias correction,
the Liquid EVM substrate, the FHE precompile activations, and the
DEX matching engine each implicate one or more patent claims.

\textbf{Remedies}: Reasonable royalty floor (35 U.S.C.\ \S 284); lost
profits where established; injunction (\S 283); treble damages for
willfulness (\S 284); attorneys' fees in exceptional cases (\S 285).

\subsection*{Count IV --- Breach of Fiduciary Duty (Del., against Jayant Sonsurkar individually)}

As a corporate officer, defendant owes the corporation duties of
care and loyalty. Mr.\ Sonsurkar's failure to identify and flag the
LEL / LRL--PR licensing mismatch throughout his tenure is at minimum
gross negligence in the CSO duty of care; if the record supports
knowledge or willful blindness, it crosses into duty-of-loyalty
breach.

\subsection*{Count V --- Aiding and Abetting Breach (Del., against Messrs.\ Choi, Tromblye, Church)}

Per \emph{Malpiede v.\ Townson}, 780 A.2d 1075 (Del.\ 2001): (i) a
fiduciary relationship; (ii) breach; (iii) knowing participation by
the non-fiduciary; (iv) damages. To the extent the named officers
had actual or constructive knowledge of the mismatch and continued to
direct commercial operations, they aided and abetted Mr.\ Sonsurkar's
breach.

\subsection*{Count VI --- Unjust Enrichment}

The Counterparty's commercial deployment of Lux IP without an
executed license enriches the Counterparty at Lux's expense; the
LEL and LRL--PR expressly negate any justification.

\subsection*{Count VII --- Conversion}

The Counterparty's deployment of LRL--PR-protected patent-bearing
modules constitutes intentional interference with Lux's exclusive
rights.

\subsection*{Count VIII --- Tortious Interference with Prospective Economic Advantage}

The Counterparty's white-label tenants (MLC, VCC, TZERO, ReTokens)
constitute particular instances of interference with Lux's
commercial-license pipeline with third-party ecosystem partners.

\subsection*{Count IX --- Securities Fraud (Rule 10b-5)}

To the extent the current \$75M / \$1.1B-valuation raise is supported by representations or
omissions regarding the Counterparty's IP ownership / licensing
status, elements of Rule 10b-5 are met. \textbf{Lux's primary
utility on this count is collateral leverage rather than direct
recovery; standing analysis is for outside counsel.}

\subsection*{Count X --- Lanham Act False Advertising (15 U.S.C.\ \S 1125(a))}

The Counterparty's marketing materials describe its technology stack
(post-quantum cryptography, FHE-based dark-pool matching, MPC
custody, multi-chain architecture, Quasar-derived consensus) in
terms that ascribe ownership or proprietary control without
disclosing the Lux origin or the conditional / unlicensed nature of
the access.

\subsection*{Count XI --- Common-Law Fraud}

To the extent Counterparty officers made representations to Lux
personnel that licensing arrangements would be formalized in due
course, and Lux relied to its detriment by continuing engineering
contribution without an executed agreement, common-law fraud lies.

\subsection*{Count XII --- Civil Conspiracy}

The Counterparty officer team's concerted commercial deployment of
unlicensed Lux IP and capital-raising on the back of that deployment
supports civil-conspiracy joint and several liability.

% ─────────────────────────────────────────────────────────────────
\section*{VI.\ Damages Analysis}

\subsection*{Small Scenario (Path A --- Negligent Posture)}

\begin{center}
\begin{tabular}{lr}
\toprule
Back-license value (12 yr at \$2.5M/yr base) & \$30M \\
Pre-judgment interest (5\%, simple) & \$10M \\
Royalty component (1.5\% of estimated GR) & \$10M \\
Forward-looking license value (NPV) & \$30M \\
\midrule
\textbf{Subtotal --- Small Scenario} & \textbf{\$80M} \\
\bottomrule
\end{tabular}
\end{center}

\subsection*{Medium Scenario (Path B --- Mixed-Knowledge, Contested)}

\begin{center}
\begin{tabular}{lr}
\toprule
Small-Scenario base & \$80M \\
DTSA 2$\times$ exemplary on willful misappropriation & +\$30M \\
Patent willfulness component (selective trebling) & +\$40M \\
Disgorgement on current-raise attribution (5\% of enterprise value, TODO: rebase per counsel) & +\$55M \\
Attorneys' fees and costs & +\$10M \\
\midrule
\textbf{Subtotal --- Medium Scenario} & \textbf{\$215M} \\
\bottomrule
\end{tabular}
\end{center}

\subsection*{Big Scenario (Path B Fully Litigated, Adverse Officer Record)}

\begin{center}
\begin{tabular}{lr}
\toprule
Base + interest + royalty & \$80M \\
DTSA full 2$\times$ exemplary & +\$80M \\
Patent full 3$\times$ trebling on willfulness & +\$120M \\
Disgorgement on current-raise attribution (15--20\% of enterprise value, TODO: rebase per counsel) & +\$200M \\
Lanham Act + Lanham fees & +\$10M \\
Securities fraud (if Lux can establish standing) & +\$50M \\
Attorneys' fees, costs, pre-judgment interest & +\$50M \\
\midrule
\textbf{Subtotal --- Big Scenario} & \textbf{\$590M} \\
\bottomrule
\end{tabular}
\end{center}

\subsection*{Opportunity-Cost Damages: 90-Day Bandwidth Diversion}

\textbf{The damages framework above does not yet capture a distinct
and significant category of harm: the opportunity cost to Lux of
having Mr.\ Kelling's --- and the broader Lux engineering team's ---
bandwidth diverted to the Counterparty's platform and to a
constellation of related ventures during a critical 90-day window.
The git record is unambiguous.}

\textbf{Quantification (rolling 90-day window) --- Counterparty
primary codebase.} Audit of the \texttt{liquidityio/*} GitHub
organization for commits attributable to Mr.\ Kelling:

\begin{center}
\begin{small}
\begin{tabular}{lrrr}
\toprule
\textbf{Repo} & \textbf{Zach commits} & \textbf{Total commits} & \textbf{Zach \%} \\
\midrule
\texttt{liquidityio/ats}        & 806    & 1{,}426 & 57\% \\
\texttt{liquidityio/bd}         & 388    & 751     & 52\% \\
\texttt{liquidityio/ta}         & 208    & 364     & 57\% \\
\texttt{liquidityio/mpc}        & 31     & 95      & 33\% \\
\texttt{liquidityio/evm}        & 23     & 32      & 72\% \\
\texttt{liquidityio/fhe}        & 10     & 14      & 71\% \\
\texttt{liquidityio/iam}        & 27     & 29      & 93\% \\
\texttt{liquidityio/node}       & 338    & 380     & 89\% \\
\texttt{liquidityio/operator}   & 354    & 406     & 87\% \\
\texttt{liquidityio/universe}   & 1{,}436 & 1{,}920 & 75\% \\
\texttt{liquidityio/contracts}  & 72     & 176     & 41\% \\
\texttt{liquidityio/dex}        & 24     & 32      & 75\% \\
\texttt{liquidityio/platform}   & 36     & 51      & 71\% \\
\texttt{liquidityio/superadmin} & 243    & 338     & 72\% \\
\texttt{liquidityio/exchange}   & 522    & 1{,}170 & 45\% \\
\texttt{liquidityio/console}    & 28     & 38      & 74\% \\
\texttt{liquidityio/ui}         & 27     & 143     & 19\% \\
\texttt{liquidityio/sdk}        & 30     & 30      & 100\% \\
\texttt{liquidityio/sdk-go}     & 9      & 9       & 100\% \\
\texttt{liquidityio/cli}        & 114    & 117     & 97\% \\
\midrule
\textbf{Subtotal (Counterparty primary, 90-day)} & \textbf{4{,}726} & \textbf{7{,}521} & \textbf{63\%} \\
\bottomrule
\end{tabular}
\end{small}
\end{center}

\textbf{Quantification --- related-venture codebases.} The audit
extends across five additional venture trees on which Mr.\ Kelling's
bandwidth was concurrently directed during the same 90-day window:

\begin{center}
\begin{small}
\begin{tabular}{lrrl}
\toprule
\textbf{Venture tree} & \textbf{Zach commits} & \textbf{LOC (add+del)} & \textbf{Note} \\
\midrule
VCC (vcc/exchange + contracts + wallet)        & 143 & 3{,}038{,}420 & \textbf{100\% new venture} (white-label tenant) \\
MLC (mlc/exchange + contracts + wallet)        & 153 & 3{,}044{,}377 & \textbf{100\% new venture} (white-label tenant) \\
EquityTable (loan-funding-engine + app)        & 98  & 681{,}666     & \textbf{100\% new venture} \\
OnyxPlus (admin + verify + onyxd + onboarding) & 188 & 57{,}379      & \textbf{Largely 100\% new venture} (KYC/ID JV) \\
LiquidityArchive (legacy + simplici-*)         & 683 & 9{,}541{,}566 & Legacy / archive maintenance \\
\midrule
\textbf{Subtotal (related ventures, 90-day)}   & \textbf{1{,}265} & \textbf{16{,}363{,}408} & \\
\bottomrule
\end{tabular}
\end{small}
\end{center}

\textbf{Grand totals across all six venture trees (90-day window):}

\begin{center}
\begin{tabular}{lr}
\toprule
\textbf{Total Zach commits (all 6 trees)} & \textbf{5{,}991} \\
\textbf{Total LOC churn (add + del across all 6 trees)} & \textbf{$\sim$93 million} \\
\textbf{LOC churn in Counterparty primary (top 8 repos)} & \textbf{$\sim$77 million} \\
\textbf{Distinct working days in \texttt{ats/} alone} & \textbf{55 (six days/week)} \\
\textbf{Repos where Zach $>$70\% of all commits} & \textbf{12 of 20 Counterparty repos} \\
\textbf{Repos where Zach $=$ 100\% of all commits} & \textbf{2 (\texttt{sdk}, \texttt{sdk-go})} \\
\bottomrule
\end{tabular}
\end{center}

\textbf{Three of the five related ventures (VCC, MLC, EquityTable) were
authored entirely from scratch during the 90-day window} ---
representing approximately \textbf{6.8 million lines} of net new code
across companies that the Counterparty offers to its tenants and
partners. A fourth (OnyxPlus, the KYC/ID joint venture) was largely
new in the same window. This is not maintenance contribution; this
is greenfield construction of multiple companies.

\textbf{Mr.\ Kelling's contribution to Lux upstream in the same
90-day window (the foundational Lux IP that powers the
Counterparty's stack):}

\begin{center}
\begin{small}
\begin{tabular}{lrr}
\toprule
\textbf{Lux repo} & \textbf{Zach commits} & \textbf{Total commits} \\
\midrule
\texttt{luxfi/node}             & \textbf{0} & 179 \\
\texttt{luxfi/consensus}        & \textbf{0} & 104 \\
\texttt{luxfi/crypto}           & \textbf{0} & 64 \\
\texttt{luxfi/evm}              & \textbf{0} & 99 \\
\texttt{luxfi/vm}               & \textbf{0} & 27 \\
\texttt{luxfi/precompile}       & \textbf{0} & 98 \\
\midrule
\textbf{Total (90-day window)} & \textbf{0} & \textbf{571} \\
\bottomrule
\end{tabular}
\end{small}
\end{center}

\textbf{The contrast is decisive}: 5,991 commits and $\sim$93M LOC
of Zach's bandwidth into Counterparty-side ventures versus exactly
zero commits to the upstream Lux repositories he is Founder and
Chief Cryptographer of.

\textbf{Value-creation rate translation.} The Counterparty's reported
current-round valuation in the active raise is approximately \$1.1 billion. The platform supporting that raise was substantially
built or rebuilt by Mr.\ Kelling during the 90-day window. The
implied value-creation rate is approximately \textbf{\$1.1B / 90
days = \$12.2M per day} of diverted attention --- and that figure
captures only the Counterparty's primary codebase; it does not yet
include the enterprise value of the three additional 100\%-new
ventures (VCC, MLC, EquityTable) and the largely-new OnyxPlus JV.
Counting those, the per-day value-creation rate of Mr.\ Kelling's
bandwidth during the 90-day window is materially higher than the
\$12.2M figure above.

\textbf{The opportunity cost to Lux} --- measured by what Lux's own
roadmap would have produced during the same window had Mr.\
Kelling's attention been available to Lux upstream --- \textbf{is
bounded above by that rate and below by the proportional founder-
share of Lux's enterprise trajectory}.

\begin{center}
\begin{tabular}{lr}
\toprule
\textbf{Component} & \textbf{Range} \\
\midrule
Bandwidth diversion at 50\% rate (90 days) & \$500M -- \$600M \\
At 25\% rate & \$250M -- \$300M \\
At 10\% rate (conservative) & \$100M -- \$120M \\
\midrule
\textbf{Conservative opportunity-cost claim} & \textbf{\$100--300M} \\
\bottomrule
\end{tabular}
\end{center}

This component is independent of and additive to the IP damages
above. \textbf{It is the single most striking quantitative fact in
the memorandum} and should be lead-in evidence for Mr.\ Dupont in
opening negotiation.

\subsection*{Damages Summary (Including Opportunity Cost)}

\begin{center}
\begin{tabular}{lrr}
\toprule
\textbf{Scenario} & \textbf{IP damages} & \textbf{+ Opp.\ Cost} \\
\midrule
Small (Path A, settlement context) & \$80M & \$180--380M \\
Medium (Path B, mixed-record litigation) & \$215M & \$315--515M \\
Big (Path B, full-record litigation w/ multipliers) & \$590M & \$690--890M \\
\bottomrule
\end{tabular}
\end{center}

% ─────────────────────────────────────────────────────────────────
\section*{VII.\ Trading-Fee Revenue Projections (User-Scale Analysis)}

The following projects the Counterparty's enterprise-value trajectory
at four user-base milestones and computes Lux's economic interest at
each under a proposed 10\% equity stake.

\subsection*{Annual Revenue, Enterprise Value, and Lux 10\% Stake}

\begin{center}
\begin{small}
\begin{tabular}{lrrrr}
\toprule
\textbf{Scenario} & \textbf{Users} & \textbf{Revenue} & \textbf{Enterprise Value} & \textbf{Lux 10\%} \\
\midrule
\multicolumn{5}{l}{\textit{Conservative ARPU (\$50/yr)}} \\
Stage 1 & 100K & \$5M & \$50M (10$\times$) & \$5M \\
Stage 2 & 1M & \$50M & \$500M (10$\times$) & \$50M \\
Stage 3 & 10M & \$500M & \$5B (10$\times$) & \$500M \\
Stage 4 & 100M & \$5B & \$50B (10$\times$) & \$5B \\
\midrule
\multicolumn{5}{l}{\textit{Base ARPU (\$100/yr)}} \\
Stage 1 & 100K & \$10M & \$150M (15$\times$) & \$15M \\
Stage 2 & 1M & \$100M & \$1.5B (15$\times$) & \$150M \\
Stage 3 & 10M & \$1B & \$15B (15$\times$) & \$1.5B \\
Stage 4 & 100M & \$10B & \$150B (15$\times$) & \textbf{\$15B} \\
\midrule
\multicolumn{5}{l}{\textit{Premium ARPU (\$200/yr)}} \\
Stage 1 & 100K & \$20M & \$400M (20$\times$) & \$40M \\
Stage 2 & 1M & \$200M & \$4B (20$\times$) & \$400M \\
Stage 3 & 10M & \$2B & \$40B (20$\times$) & \$4B \\
Stage 4 & 100M & \$20B & \$400B (20$\times$) & \textbf{\$40B} \\
\bottomrule
\end{tabular}
\end{small}
\end{center}

\textbf{Lux's 10\% equity, retained through the growth curve, is
worth between \$5M (Stage 1 Conservative) and \$40B (Stage 4
Premium).} Mid-case Base ARPU at 100M users implies \$15B of Lux
equity value.

% ─────────────────────────────────────────────────────────────────
\section*{VIII.\ Strategic Partnership Upside}

Lux brings substantial commercial value to a partnership-resolution
of this dispute. This section enumerates what the deal team can
offer as upside in any negotiation tier.

\subsection*{First Major Client: AvaTrade --- \$70B/month Trading Volume}

\textbf{Lux has secured AvaTrade as the first major client for the
combined Lux + Counterparty venue, with an estimated \$70 billion
per month in trading volume routed through the joint architecture.}
AvaTrade is an established multi-asset broker with operations across
multiple regulated jurisdictions; the partnership agreement
underlying this arrangement is documented separately at
\texttt{\textasciitilde/work/lux/legal/partnerships/avatrade/Lux\_AvaTrade\_Strategic\_Partnership\_Agreement.pdf}.

At 8 bps blended take-rate on \$70B/mo = \$5.6M/mo = \textbf{\$67M
per year in fee revenue from this single client alone}. At Stage-2-
plus volume routing, the AvaTrade contribution materially exceeds
the Counterparty's current revenue base and is, by itself, the
strongest argument for partnership over litigation.

\subsection*{Crypto-Exchange Conversion at 8 Basis Points (Floor)}

A specific deliverable Lux brings under the partnership: \textbf{the
ability to offer crypto-exchange conversion at 8 basis points} ---
materially undercutting the incumbent fee structure at major US
crypto venues (Coinbase / Kraken / Gemini at 50--200 bps; Robinhood
PFOF-equivalent at 35--85 bps; institutional desks at 5--15 bps
floor but with size minimums). The 8-bps floor is enabled by Lux's
direct exchange and liquidity infrastructure and is a competitive
moat for the joint Lux + Counterparty venue.

\subsection*{OnyxPlus / Telos JV ID Verification at Internal Cost}

The Lux--Telos joint venture (OnyxPlus) provides identity verification
infrastructure operating at the Lux ecosystem's internal cost basis
--- materially below merchant pricing for equivalent
KYC/AML/document/biometric services. Under the SCLA, \textbf{the
Counterparty receives access to OnyxPlus ID verification at Lux's
internal cost, not at market rates}. The annual savings depend on
the Counterparty's user-onboarding velocity but at Stage-2-plus
scale run into the tens of millions per year.

\subsection*{Co-Branding for Lux-Network Chains}

\textbf{All Lux-network L1/L2/L3/L4 chains receive explicit
co-branding on the Counterparty's customer-facing surfaces}: the
Counterparty's marketing, the venue's chain selector, the
SecurityToken issuance pages, and any institutional partnership
collateral. This brand co-positioning is a non-negotiable term
under the SCLA (\S XI).

\subsection*{Real-World Asset Pipeline}

Lux operates the infrastructure layer behind multi-billion-dollar
RWA tokenization pipelines. Post-SCLA, the Counterparty is the
natural US-regulated listing destination:

\begin{itemize}[leftmargin=1.5em,itemsep=0.1em]
\item Tokenized US securities + private equity (Carta ingestion).
\item Tokenized real estate / NNN-SPV-wrapped property exposures.
\item Tokenized commodities (gold, silver, FX, energy).
\item Tokenized fund LP interests.
\item Tokenized fixed-income.
\item Tokenized music + IP royalty streams.
\item Tokenized art and collectibles ($>$\$1M lots).
\end{itemize}

\textbf{Conservative estimate of incremental deal flow under the
SCLA: \$5--15 billion in nominal asset value over the first 12--24
months post-close.}

\subsection*{Reciprocal License to Counterparty ATS / BD / TA / Listed Securities --- Lux's Forward Venue Strategy}

\textbf{Lux Industries Inc.\ does not presently operate an
alternative trading system or broker-dealer in the United States.}
However, Lux's historical regulatory posture and its forward strategy
include venue licenses in multiple jurisdictions, and the reciprocal
license to the Counterparty's ATS / BD / TA / listed-securities
surfaces described in \S XI(9) is intended to give Lux operational
optionality across that footprint as it matures.

\textbf{Historical native support:} Lux (and Lux Partners Limited
as predecessor) previously operated regulatory arrangements in the
\textbf{Isle of Man (IOM)} (IOM Financial Services Authority
supervision) and \textbf{Luxembourg} (Commission de Surveillance du
Secteur Financier engagement). The IOM and Luxembourg footprints
provided native EU-passportable and offshore-fund-eligible regulatory
reach.

\textbf{Forward expansion:} Lux is actively evaluating venue-license
applications across the United States (Reg ATS), United Kingdom
(FCA MTF or equivalent), European Union (MiCA-compliant Trading
Platform), Singapore (MAS Capital Markets Services Licence +
Recognised Market Operator), and the UAE (VARA). Each license,
when obtained, becomes a venue at which Lux may operate the
Counterparty-licensed ATS / BD / TA technology under \S XI(9).

\textbf{Cross-listing rights:} Every security that lists on the
Counterparty venue is available for cross-listing on any Lux-network
or Lux-affiliate venue, and every security that lists on a Lux
venue is available for cross-listing on the Counterparty's venue.

\subsection*{Banking Licenses --- SF Private Bank Partnership and Global Coverage}

\textbf{Lux holds a strategic partnership with SF Private Bank
(``SFPB'') and adjacent financial institutions that, in combination,
provide regulated coverage on every populated continent.} The
combined licensed footprint:

\begin{itemize}[leftmargin=1.5em,itemsep=0.1em]
\item \textbf{North America}: United States (federal banking partner,
FDIC-insured rails); Canada (chartered banking partner, CDIC-insured).
\item \textbf{Europe}: Sweden (Finansinspektionen-supervised, EU
passport); historical IOM (Financial Services Authority); historical
Luxembourg (CSSF); ongoing UK (FCA) and Switzerland (FINMA) expansion.
\item \textbf{Asia-Pacific}: Singapore (MAS Capital Markets Services
and banking); New Zealand (RBNZ-licensed deposit-taker); ongoing
Hong Kong (SFC type-1 / type-7) and Japan (JVCEA) expansion.
\item \textbf{Middle East}: UAE (VARA-supervised; DIFC / ADGM
counterparty relationships).
\item \textbf{Latin America}: regional partnerships under bilateral
arrangements; specific licensing on a case-by-case basis.
\item \textbf{Africa}: pilot arrangements; specific licensing on a
case-by-case basis.
\end{itemize}

Services Lux can deliver under the SCLA include Banking-as-a-Service
(BaaS), Payment Service Provider (PSP) rails, and \textbf{a custody
solution that materially undercuts the Fireblocks and BitGo
incumbents on price} (built on the same MPC infrastructure already
in the Counterparty's production).

\textbf{Estimated independent cost for the Counterparty to acquire
equivalent global license suite: \$57--140M in direct regulatory
and capital expenditure, plus 5--7 years of timeline.} See Schedule
C.

\subsection*{International Distribution: M7 and 1B+ End-User Reach}

Through Hanzo Inc., the broader corporate family, and the Lux
ecosystem partnerships, the M7 global initiatives provide
distribution access to more than 1 billion end-users across consumer,
institutional, and retail channels. Post-SCLA, Lux's international
counterparty network operates as a Reg-S-compatible distribution
channel for Counterparty-listed assets.

\subsection*{Hanzo Inc.\ Portfolio Precedents}

Hanzo Inc., the venture studio under which the Lux relationship has
historically operated, has applied the IP-in-kind founder-equity
model successfully to multiple spawned unicorn-class companies:

\begin{itemize}[leftmargin=1.5em,itemsep=0.1em]
\item \textbf{Damon Motorcycles} --- electric / smart-motorcycle company.
\item \textbf{Triller} --- short-form video / social platform.
\item \textbf{Bellabeat} --- women's health wearables.
\item \textbf{Unicorn Gold} --- precious-metals tokenization platform.
\item \textbf{Lux Network} --- itself a Hanzo-spawn outcome.
\item \textbf{Zoo Labs Foundation / Zoo Network} --- decentralized AI
research network.
\end{itemize}

Retained-parent-equity range across the portfolio: 10\%--40\%. The
10\% ask in the Counterparty is at the bottom of that internal range.

\subsection*{The \$100B--\$1T Target Institution}

The combined Lux + Counterparty vision is the construction of a
\$100-billion-to-\$1-trillion-class institution at the intersection
of tokenized real-world assets, post-quantum-secure infrastructure,
regulated US securities markets, and global distribution. The deal
team should treat the negotiation not as a zero-sum settlement but
as the founding-document of a partnership creating value materially
in excess of the negotiating range.

% ─────────────────────────────────────────────────────────────────
\section*{IX.\ Wait vs.\ Settle Analysis}

\textbf{Forces favoring delay:} equity appreciation; royalty accrual;
current-raise disgorgement expansion; pattern-of-conduct evidence.

\textbf{Forces favoring settlement now:} statutory damages caps don't
grow with delay (DTSA at 2$\times$; patent at 3$\times$); dilution risk
on Lux's equity component; operational risk (engineer-around); evidence
preservation; DTSA 3-year limitations from discovery; Lux operating
runway; negotiating leverage compounds short-term then evaporates if
the Counterparty hits trouble; public visibility window.

\textbf{Recommendation:} \textbf{Settle now, with strong protective
terms} --- significant cash component, 10\%+ equity with weighted-
average broad-based anti-dilution, full ratchet on down-rounds, pro-
rata participation, tag-along and co-sale rights, 1+1 mandatory board
seats, 99-year exclusive USA-market license to the Counterparty
(conditional), patent grant, AvaTrade routing, 8 bps conversion floor,
OnyxPlus access at internal cost, co-branding for Lux chains, and
reciprocal permanent license to Counterparty ATS / BD / TA / listed-
securities / cross-listings (\S XI). \textbf{Statutory damages caps
do not grow with delay; the equity does.} Locking equity now with
strong anti-dilution captures the future-growth narrative without
the operational risks of continued reliance on a deteriorating
informal arrangement.

% ─────────────────────────────────────────────────────────────────
\section*{X.\ Settlement Triangulation (Cash / Equity / Features Matrix)}

\subsection*{Settlement Sizing Tiers}

\begin{center}
\begin{small}
\begin{tabular}{lrrrr}
\toprule
\textbf{Tier} & \textbf{Cash} & \textbf{Equity} & \textbf{Lux Nominal} & \textbf{Implied Posture} \\
\midrule
Floor (S--)       & \$5M   & 5\%   & \$60M   & Soft settle; preserve relationship \\
Small (S)         & \$10M  & 10\%  & \$120M  & Approach baseline (recommended floor) \\
Medium (M)        & \$20M  & 15\%  & \$185M  & Path B leverage applied softly \\
Big (B)           & \$30M  & 25\%  & \$305M  & Full Path B post-investigation \\
Aspirational (A+) & \$50M  & 35\%  & \$435M  & Litigation-credible ceiling \\
50/50 (max)       & \$100M & 50\%  & \$650M  & Full partnership re-cap \\
\bottomrule
\end{tabular}
\end{small}
\end{center}

\subsection*{Cash vs.\ Equity Mix Sensitivity (at \$120M floor)}

\begin{center}
\begin{tabular}{lll}
\toprule
\textbf{Mix} & \textbf{Cash} & \textbf{Equity at \$1.1B valuation} \\
\midrule
All-cash & \$120M & 0\% \\
Heavy cash & \$60M & 5.5\% \\
Balanced (recommended) & \$10M & 10\% \\
Light cash & \$5M & 10.5\% \\
All-equity & \$0 & 11\% \\
\bottomrule
\end{tabular}
\end{center}

\subsection*{Feature-Unlock by Settlement Tier}

\begin{center}
\begin{small}
\begin{tabular}{p{6cm} ccccc}
\toprule
\textbf{Forward-value feature} & \textbf{S--} & \textbf{S} & \textbf{M} & \textbf{B} & \textbf{A+} \\
\midrule
LEL/LRL--PR perpetual license & $\checkmark$ & $\checkmark$ & $\checkmark$ & $\checkmark$ & $\checkmark$ \\
Patent grant on covered IP & $\checkmark$ & $\checkmark$ & $\checkmark$ & $\checkmark$ & $\checkmark$ \\
Maintenance \& support of upstream OSS & & $\checkmark$ & $\checkmark$ & $\checkmark$ & $\checkmark$ \\
99-year exclusive USA-market license & & $\checkmark$ & $\checkmark$ & $\checkmark$ & $\checkmark$ \\
1 Lux + 1 independent director seat & & $\checkmark$ & $\checkmark$ & $\checkmark$ & $\checkmark$ \\
RWA pipeline routing into Counterparty & & $\checkmark$ & $\checkmark$ & $\checkmark$ & $\checkmark$ \\
AvaTrade flow routing into venue & & $\checkmark$ & $\checkmark$ & $\checkmark$ & $\checkmark$ \\
8 bps crypto-conversion floor & & $\checkmark$ & $\checkmark$ & $\checkmark$ & $\checkmark$ \\
OnyxPlus ID verify at internal cost & & $\checkmark$ & $\checkmark$ & $\checkmark$ & $\checkmark$ \\
Co-branding for Lux-network chains & & $\checkmark$ & $\checkmark$ & $\checkmark$ & $\checkmark$ \\
Reciprocal license to ATS/BD/TA & & & $\checkmark$ & $\checkmark$ & $\checkmark$ \\
SFPB banking-license access (BaaS, PSP) & & & $\checkmark$ & $\checkmark$ & $\checkmark$ \\
Custody undercutting Fireblocks & & & $\checkmark$ & $\checkmark$ & $\checkmark$ \\
International M7 distribution channels & & & & $\checkmark$ & $\checkmark$ \\
Closed-source (\texttt{luxcpp}) GPU access & & & & & $\checkmark$ \\
Joint go-to-market programs & & & & $\checkmark$ & $\checkmark$ \\
\bottomrule
\end{tabular}
\end{small}
\end{center}

\subsection*{Damages-vs-Settlement Triangulation}

\begin{center}
\begin{tabular}{lrr}
\toprule
\textbf{Outcome path} & \textbf{Lux receives} & \textbf{Counterparty pays} \\
\midrule
Settle at Floor (S--) & \$60M nominal & \$5M cash + 5\% equity \\
Settle at Small (S, recommended) & \$120M nominal & \$10M cash + 10\% equity \\
Settle at Medium (M) & \$185M nominal & \$20M cash + 15\% equity \\
Settle at Big (B) & \$305M nominal & \$30M cash + 25\% equity \\
Settle at Aspirational (A+) & \$435M nominal & \$50M cash + 35\% equity \\
Litigate: Small (Path A) & \$80M & \$80M cash, no equity \\
Litigate: Medium (Path B mixed) & \$215M & \$215M, no equity \\
Litigate: Big (Path B full) & \$590M & \$590M, no equity \\
Litigate + Opp.\ Cost (Big) & \$690--890M & idem, no equity, no relationship \\
\bottomrule
\end{tabular}
\end{center}

\textbf{Settlement dominates litigation at every tier for both sides,
provided the equity component is properly valued.}

% ─────────────────────────────────────────────────────────────────
\section*{XI.\ Mandatory Terms (Lux's Floor in Any Settlement)}

The following terms are non-negotiable. Any settlement at any tier
must include each of them.

\textbf{1.} \textbf{Board composition --- expansion to a minimum of
seven (7) directors total, with structured composition.} The
Counterparty's existing three-director board failed to identify the
licensing mismatch over the relevant period despite having a
maintainer-access-empowered CSO whose duty included precisely that
review. \textbf{The 3-person board structure is, in light of that
record, an inadequate governance vehicle going forward}, and the
disinterested-investor case for board expansion is unambiguous.

The mandatory composition is:

\begin{itemize}[leftmargin=2em,itemsep=0.1em]
\item \textbf{Three (3) existing director seats retained} from the
current board.
\item \textbf{One (1) Lux director seat}, for Lux Industries Inc.'s
direct representative (Mr.\ Pahlavi's designate).
\item \textbf{Three (3) mutually-agreed independent director seats},
each nominee with no prior commercial or personal relationship to
either Lux or the Counterparty, serving as the ongoing governance
anchor of the SCLA.
\end{itemize}

Effective at closing. \textbf{No ownership-maintenance condition
attaches to the Lux seat.} The structure ensures (a) no single
faction (existing officers, Lux, or independents) holds majority
voting power; (b) the independent block of three directors provides
the disinterested-quorum threshold for any related-party transaction
under Delaware General Corporation Law \S 144; and (c) the
Counterparty's institutional governance moves from informal,
founder-controlled posture to formal, multi-faction supervision
appropriate to a venue valued in the \$1B+ range with FINRA / SEC
exam exposure.

In the alternative (a 7-director soft floor that Lux will accept
under appropriate compensating concessions), the composition may
shift to: 2 existing seats retained + 2 Lux directors + 3 independent
seats. This variant displaces one of the existing seats and
strengthens Lux's voting position; Lux's deal team may invoke this
alternative as leverage if the Counterparty resists the principal
structure above.

\textbf{2.} \textbf{99-year exclusive USA-market license to the
Counterparty} for the LEL/LRL--PR portfolio in scope. Conditional on
maintenance of all SCLA covenants. Material breach terminates
exclusivity. Outside the USA, Lux retains the right to license to
other parties.

\textbf{3.} \textbf{Anti-dilution: weighted-average broad-based,
with full ratchet on any down-round below the Lux issuance price.}

\textbf{4.} \textbf{Pro-rata participation right} on all future
Counterparty preferred financings.

\textbf{5.} \textbf{Tag-along and co-sale rights} on any founder
transfer exceeding 1\% of fully-diluted equity.

\textbf{6.} \textbf{Information rights}: quarterly and annual
financials; all board materials and minutes; all share-issuance and
option-grant resolutions; counsel and accountant access under NDA.

\textbf{7.} \textbf{Patent grant} on all current and future Lux
patents covering the licensed IP. Worldwide; perpetual; royalty-free
under the SCLA; expressly inclusive of patents not yet filed at SCLA
execution.

\textbf{8.} \textbf{Maintenance and support of upstream Lux OSS}:
continuing on the shared-engineering basis as today, no separate
maintenance fee.

\textbf{9.} \textbf{Reciprocity --- Permanent license to Counterparty
ATS, BD, TA, listed securities, and cross-listings.} As symmetric
consideration for Lux's perpetual IP license, the Counterparty grants
Lux Industries Inc.\ (and any Lux-controlled affiliate or Lux-network
L1/L2/L3/L4 chain maintaining Lux consensus and Warp interop):

\begin{itemize}[leftmargin=2em,itemsep=0.1em]
\item \textbf{Permanent, royalty-free, non-exclusive license to the
Counterparty's ATS technology and operational know-how}, sufficient
for Lux to operate (or sublicense to a Lux-controlled regulated
entity) an ATS in any jurisdiction where Lux holds or obtains the
necessary venue license.
\item \textbf{Permanent, royalty-free, non-exclusive license to the
Counterparty's BD technology and know-how} on equivalent terms.
\item \textbf{Permanent, royalty-free, non-exclusive license to the
Counterparty's TA technology and know-how} on equivalent terms.
\item \textbf{Reciprocal listing and cross-listing rights}: every
security listed on the Counterparty venue is available for
cross-listing on any Lux venue, and every security listed on a Lux
venue is available for cross-listing on the Counterparty venue, on
no-less-favorable terms.
\item \textbf{Sublicensing rights}: Lux may sublicense to Lux-
controlled jurisdictions and operating affiliates, including the
IOM, Luxembourg, FCA, MAS, VARA, MiCA, and FINMA jurisdictional
expansions described in \S VIII.
\end{itemize}

The reciprocal grant runs 99 years, conditional on continued SCLA
performance.

\textbf{10.} \textbf{Co-branding for Lux-network chains.} All Lux-
network L1/L2/L3/L4 chains receive explicit co-branding on the
Counterparty's customer-facing surfaces (marketing, venue chain
selector, SecurityToken issuance pages, and institutional partnership
collateral).

\textbf{11.} \textbf{OnyxPlus ID verification at Lux's internal
cost.} The Counterparty receives access to OnyxPlus
KYC/AML/document/biometric verification services at Lux's internal
cost basis, not at market rates.

\textbf{12.} \textbf{AvaTrade routing}: the AvaTrade flow secured by
Lux is routed through the joint venue post-SCLA; the Counterparty
receives the proportional benefit of the trading volume on terms
consistent with the venue's standard listing arrangements.

\textbf{13.} \textbf{8 bps crypto-conversion floor}: Lux's
infrastructure enables the joint venue to offer crypto-exchange
conversion at 8 basis points (floor); the Counterparty's existing
crypto-conversion pricing is repriced to this floor under the SCLA.

\textbf{14.} \textbf{Engineering commitment}: the Counterparty
activates the Lux Warp precompile (canonical address
\texttt{0x0200000000000000000000000000000000000005}) on every Liquid
chain at the next chain-upgrade window. New Liquid chains include
the \texttt{warpConfig} block in genesis.

\textbf{15.} \textbf{Litigation hold and document preservation}: the
Counterparty agrees to preserve all records, communications, and
code-review artifacts relating to the licensing posture from 2014 to
the SCLA execution date.

\textbf{16.} \textbf{Independent investigation cooperation}: the
Counterparty cooperates fully with any independent investigation of
officer conduct undertaken by the disinterested directors or their
appointed counsel.

% ─────────────────────────────────────────────────────────────────
\section*{XII.\ Special Note Re: Mr.\ Zach Kelling}

\textbf{Mr.\ Kelling occupies a position requiring particular care
in Lux's litigation and settlement strategy.} The factual context
established in \S IV(F) bears directly here: Mr.\ Kelling was the
\emph{remediator} of the pre-Kelling Counterparty technical and
security failures (the broken time-price priority matching; the
\$40,000 ACH-exploit theft) and the architect of the rebuilt platform
that supports the current \$75M / \$1.1B-valuation raise. He was not the author
of the conduct giving rise to this matter; he is, if anything, the
party whose work was misappropriated alongside Lux's.

\textbf{Litigation-strategy implications:}

\begin{itemize}[leftmargin=1.5em,itemsep=0.1em]
\item Mr.\ Kelling should \emph{not} be named as a defendant in any
Lux suit. The conduct alleged is not his.
\item Mr.\ Kelling is likely a material witness for Lux on (a) the
authorship and ownership of the Lux IP, (b) the absence of any
executed commercial license, and (c) the operational-displacement
narrative explaining why the mismatch was not raised internally.
\item Counsel should consider a litigation-cooperation agreement,
witness-protection arrangement, or formal indemnification of Mr.\
Kelling (as a non-defendant witness) to ensure his continued
availability without undue personal exposure.
\item In any public communication, Lux should take care not to
attribute the misconduct to Mr.\ Kelling. The narrative should be
specific to the Counterparty officers actually responsible.
\item In the SCLA itself, Lux should provide for Mr.\ Kelling's
protective treatment, including (a) a release of any claims Lux
might assert against him personally, (b) acknowledgment of his
historical authorship, and (c) a ROFR / similar mechanism if his
Counterparty role generates IP that should flow back to Lux upstream.
\end{itemize}

\textbf{Mr.\ Kelling can serve as a productive intermediary in the
Lux-Counterparty negotiation}, given his founder status at both
entities, his historical relationships, and his absence of financial
interest in the Counterparty. Mr.\ Dupont should consider engaging
Mr.\ Kelling directly in informational meetings during the
negotiation phase, subject to outside counsel's oversight.

% ─────────────────────────────────────────────────────────────────
\section*{XIII.\ Procedural Recommendations}

\textbf{Decision chain:}

\begin{enumerate}[leftmargin=1.5em,itemsep=0.1em]
\item \textbf{Mr.\ Pahlavi (President \& CEO)} --- final approval
on any executed SCLA, litigation filing, or strategic public
communication.
\item \textbf{Mr.\ Dupont (Managing Partner, Strategic Transactions;
largest investor)} --- deal lead. Conducts negotiation; selects
settlement tier; presents executed term sheet to Mr.\ Pahlavi.
\item \textbf{Mr.\ Donohue (General Counsel)} --- legal lead.
Drafts SCLA; supervises outside counsel; manages litigation hold;
leads pre-suit investigation; prepares SCLA for execution.
\item \textbf{Mr.\ Bin (CTO; memorandum author)} --- operational and
technical support. Provides factual development on IP portfolio,
licensing record, and Counterparty-side operational context.
\end{enumerate}

\textbf{Recommended outside engagement:}

\begin{itemize}[leftmargin=1.5em,itemsep=0.1em]
\item \textbf{Outside litigation counsel}: a firm with DTSA + patent
+ Delaware-fiduciary-duty experience (Wilson Sonsini, Cooley, Fenwick
\& West, or Quinn Emanuel).
\item \textbf{Valuation firm}: independent fairness opinion on
equity-for-IP allocation (Houlihan Lokey, Duff \& Phelps / Kroll, or
specialist boutique).
\item \textbf{Investigative firm}: separate firm to investigate
Counterparty officer conduct (Kroll Investigations, FTI Consulting,
or specialist tech-IP investigators).
\end{itemize}

\textbf{Target close: 75 days} from Mr.\ Pahlavi's approval to
authorize the deal team, per the milestone schedule:

\begin{center}
\begin{tabular}{lp{10cm}}
\toprule
\textbf{Day} & \textbf{Action} \\
\midrule
0   & Mr.\ Pahlavi approves; outside counsel engaged. \\
7   & Litigation hold notice transmitted to Counterparty + named officers. \\
14  & Investigation firm begins document review. \\
21  & Term sheet delivered to Counterparty at selected tier. \\
30  & Counterparty disinterested-director response received. \\
45  & Investigation report finalized; settlement terms tightened/expanded. \\
60  & Definitive SCLA executed. \\
75  & Closing: cash payment, share issuance, license activation. \\
\bottomrule
\end{tabular}
\end{center}

% ─────────────────────────────────────────────────────────────────
\section*{XIV.\ Prayer for Relief (Template)}

In any litigation, the complaint's prayer will request:

\begin{enumerate}[leftmargin=1.5em,itemsep=0.1em]
\item Compensatory damages on each Count, anticipated \$80M to \$590M
on IP claims; \$180M to \$890M with opportunity-cost component.
\item Disgorgement of unjust enrichment from the current and any future raises (rebased to enterprise value, not round size).
\item Reasonable royalty (35 U.S.C.\ \S 284; 18 U.S.C.\ \S 1836(b)(3)(B)).
\item Exemplary damages 2$\times$ on DTSA (18 U.S.C.\ \S 1836(b)(3)(C)).
\item Treble damages on patent (35 U.S.C.\ \S 284).
\item Statutory damages on copyright (17 U.S.C.\ \S 504(c)) at elected
maximum.
\item Attorneys' fees (18 U.S.C.\ \S 1836(b)(3)(D); 17 U.S.C.\ \S 505;
35 U.S.C.\ \S 285; 15 U.S.C.\ \S 1117(a)).
\item Preliminary and permanent injunctive relief enjoining further
use absent executed commercial license.
\item Pre- and post-judgment interest.
\item Such other and further relief as the Court deems just.
\end{enumerate}

\textbf{Demand for jury trial} on all triable issues.

% ─────────────────────────────────────────────────────────────────
\vspace{1em}
\hrule height 0.4pt
\vspace{0.5em}
\begin{center}\textbf{\textsc{Schedules}}\end{center}
\vspace{0.5em}

\section*{Schedule A --- Lux IP Inventory in Production}

The full inventory is maintained at the companion license audit. Headline:

\begin{itemize}[leftmargin=1.5em,itemsep=0.1em]
\item BSD-3 / Apache-2.0 / MIT modules: approximately 25 (no Lux
action required).
\item LEL modules: \texttt{aml}, \texttt{broker}, \texttt{captable},
\texttt{compliance}, \texttt{corona}, \texttt{crypto}, \texttt{fhe},
\texttt{hsm}, \texttt{maker}, \texttt{mpc}, \texttt{precompile},
\texttt{ringtail}, \texttt{threshold}, \texttt{transfer},
\texttt{treasury}, \texttt{warp}.
\item LRL--PR modules: \texttt{vm}, \texttt{staking}, \texttt{sampler},
\texttt{dex} (Go bindings).
\item LGPL modules: \texttt{evm}.
\item Modules without LICENSE: \texttt{cex} and small number of
related; counsel should treat as no-rights-granted pending upstream
remediation.
\end{itemize}

\section*{Schedule B --- Lux IP Pipeline (Available Under SCLA)}

\begin{itemize}[leftmargin=1.5em,itemsep=0.1em]
\item AI inference and compute precompiles (LP-9080 family).
\item Advanced ZK rollup substrates.
\item Post-quantum hybrid signature schemes (X-Wing + ML-KEM; hybrid
ML-DSA + ECDSA).
\item Threshold MPC dynamic re-sharing protocols.
\item Confidential-execution attestation (TEE / SEV-SNP / TDX binding
to consensus).
\item White-label partner portal infrastructure.
\item Closed \texttt{\textasciitilde/work/luxcpp/*} (separate
Private-Moat Addendum required).
\end{itemize}

\section*{Schedule C --- Banking License Cost-to-Replicate}

\begin{center}
\begin{tabular}{lr}
\toprule
\textbf{License} & \textbf{Estimated Cost} \\
\midrule
US national or state banking charter & \$20--50M \\
Canadian Schedule II / partnership & \$10--20M \\
MAS Capital Markets Services + banking (Singapore) & \$5--10M \\
RBNZ deposit-taker (New Zealand) & \$5--15M \\
Sweden Finansinspektionen (EU passport) & \$5--10M \\
UK FCA EMI / Payment Institution & \$2--5M \\
Switzerland FINMA digital-asset banking & \$5--15M \\
Hong Kong SFC Type 1 / Type 7 + VATP & \$5--15M \\
\midrule
\textbf{Total direct regulatory + capital} & \textbf{\$57--140M} \\
Plus 5--7 years of timeline + execution risk & \\
\bottomrule
\end{tabular}
\end{center}

\section*{Schedule D --- Hanzo Inc.\ Portfolio Comparables}

\begin{itemize}[leftmargin=1.5em,itemsep=0.1em]
\item Damon Motorcycles --- electric / smart-motorcycle company.
\item Triller --- short-form video / social platform.
\item Bellabeat --- women's health wearables.
\item Unicorn Gold --- precious-metals tokenization platform.
\item Lux Network --- subject of this memorandum.
\item Zoo Labs Foundation / Zoo Network --- decentralized AI research network.
\end{itemize}

Retained-parent-equity range: 10\%--40\%. The 10\% ask sits at the
floor of that range.

\section*{Schedule E --- Mandatory Terms (Lux's Floor)}

See \S XI above. Sixteen mandatory terms, summarized:

\begin{enumerate}[leftmargin=1.5em,itemsep=0.05em]
\item 1 Lux director + 1 independent director seat (closing).
\item 99-year exclusive USA-market license.
\item Weighted-average broad-based anti-dilution + full ratchet
on down-rounds.
\item Pro-rata participation rights.
\item Tag-along and co-sale rights.
\item Information rights.
\item Patent grant (current and future Lux patents).
\item Maintenance and support of upstream OSS.
\item Reciprocal permanent license to Counterparty ATS / BD / TA /
listed securities / cross-listings.
\item Co-branding for Lux-network chains.
\item OnyxPlus ID verification at Lux internal cost.
\item AvaTrade routing into joint venue.
\item 8 bps crypto-conversion floor.
\item Warp precompile activation at next chain-upgrade.
\item Litigation hold and document preservation.
\item Independent investigation cooperation.
\end{enumerate}

\vfill
\hrule height 0.4pt
\vspace{0.5em}
\begin{small}
\textit{End of memorandum. This document is confidential, privileged,
and prepared in anticipation of litigation as well as for board
strategy purposes. The proposed settlement framework is non-binding
until executed as a definitive Settlement and Cross-License Agreement.
Distribution outside Lux Industries Inc., its Board of Directors, the
named decision-chain officers (Messrs.\ Pahlavi, Dupont, Donohue, and
Bin), and engaged outside counsel requires the express written
authorization of Mr.\ Pahlavi and Mr.\ Donohue.}
\end{small}

\end{document}
